\chapter{El agua que baja}
\label{cap:aguabaja}

Este libro lleva diecisiete capítulos hablando de agua, y una parte entera, la cuarta, lleva ese nombre; pero todavía no lo dijo todo junto.

La cuenta es fácil de hacer y conviene hacerla antes de seguir. Las palabras \textit{agua},
\textit{caudal} y \textit{concesión} aparecen más de cuatrocientas veces en los diecisiete
capítulos anteriores. Aparecen en la expropiación de 1972, en los veinticinco actos de línea de ribera, en las canteras sobre el cauce, en los certificados que se le exigen a un loteo, en las
defensas que se construyen desde 1910, en la ordenanza municipal de 1914 y en el amparo que
todavía tramita. La cuarta parte reúne los instrumentos del agua; el agua, además, está en todos los demás capítulos.

Este capítulo agrega pocos documentos. \textbf{Sobre todo agrega la frase que los otros capítulos permiten
decir y no dicen}, y remite a ellos en lugar de repetirlos.

\section{La frase}

\textbf{La cuenca del río La Caldera produce agua para una ciudad que no es la suya, y el
departamento que la produce no tiene red.}

Dicho así parece una denuncia. No lo es, y el resto del capítulo intenta mostrar por qué: cada
uno de los actos que componen esa situación fue legal, fundado y en su momento razonable. Lo
que este libro sostiene es que \textbf{la suma tiene una dirección}, y que esa dirección nunca
fue decidida por nadie en particular.

\section{Lo que sale}

Tres obras, separadas por más de noventa años, y las tres llevan el agua en el mismo sentido.

\textbf{El embalse.} En 1972 se expropiaron \textbf{seiscientas ochenta y seis hectáreas} para
el dique Campo Alegre. El capítulo \ref{cap:expropiacion} reconstruye quiénes perdieron esa
tierra y cómo se les pagó. El acueducto que sale del embalse, inaugurado en 2021, se construyó para abastecer a la ciudad de Salta, y la captación más antigua sobre el río ya lo hacía (capítulo \ref{cap:expropiacion}).

\textbf{El Acueducto Norte.} En agosto de 2017 el secretario de Recursos Hídricos se reúne con
la contratista, la consultora y la empresa de agua para coordinar la determinación de la línea
de ribera en la zona de emplazamiento de la planta de tratamiento, y para fijar los requisitos
de los cruces de Los Naranjos, Guaranguay y Wierna. La obra tendrá \textbf{mil cien metros
cúbicos por hora} y abastecerá, dice el parte oficial, «a toda la zona norte de la ciudad de
Salta y a los pobladores de Vaqueros». El capítulo \ref{cap:ribera} lo documenta.

\textbf{Y la ley de 1929.} La Ley provincial 1269 (original 11078) autoriza a la Municipalidad del
Departamento de Campo Santo a reconcentrar el agua \textbf{desde las juntas del río Wierna con
el de La Caldera}, y a vigilar y prohibir nuevas bocas-toma. El capítulo \ref{cap:redes} la
transcribe y señala lo que revela: la cuenca es una y atraviesa varios departamentos, y el
arreglo institucional consistió en darle a uno de ellos potestad sobre el tramo que le
importaba. \textbf{No se creó una autoridad de cuenca: se repartió el río.} La ley ya no rige ---el Digesto la clasifica entre las de objeto cumplido o no vigentes---, pero la autoridad de cuenca que no creó sigue sin existir.

\section{Lo que no llega}

El capítulo \ref{cap:redes} se titula «Las redes que no están» y hace el inventario: qué
emprendimientos resolvieron su abastecimiento por cuenta propia y qué obras públicas quedaron
proyectadas.

Los dos movimientos conviven sin tocarse. \textbf{Los clubes de campo y los loteos resuelven su agua con
captación propia ---río, acequia, dren o pozo--- y concesión; el municipio no tiene con qué extender una red.} Cada solución
privada es, además, perfectamente legal: el régimen prevé que quien quiera el servicio lo
costee.

\textbf{Y tiene una figura jurídica propia, que este libro no había identificado.}

\begin{ficha}{Consorcios de agua pública del departamento}
\textbf{Consorcio de Agua Pública Loteo Valle Alegre}, expediente 0090034-86490 ---iniciado en
2013 y con cuerpos agregados en 2018---, con asambleas convocadas por la autoridad hídrica en
\textbf{2024} y \textbf{2026}.

\textbf{Consorcio de Usuarios de Agua Valle Hermoso --- Coronado Cueto}, expediente
0090034-181690/2020-0, con asambleas en \textbf{2022} y \textbf{2023}.

Las convocatorias las publica la propia Secretaría de Recursos Hídricos en el Boletín Oficial,
bajo el tipo de aviso «asambleas de entidades civiles».
\end{ficha}

Son los \textbf{consorcios de usuarios} del Título VI del Código de Aguas: personas jurídicas
de derecho público, entes públicos no estatales, que administran el agua desde una toma común
y a los que el artículo 195 incorpora obligatoriamente a todo usuario del área demarcada.

\begin{cautela}
\textbf{Esto matiza el apartado anterior y conviene que quede dicho.} La solución privada de
los loteos no es un pozo sin marco: es \textbf{una figura del propio Código de Aguas},
constituida por acto de la autoridad, con estatuto aprobado, asambleas publicadas y contralor
del organismo.

\textbf{Lo que no cambia es la dirección.} Dos consorcios de agua para consumo, los dos de urbanizaciones, ninguno del casco urbano ni de los parajes; los otros que el expediente del amparo registra ---el de usuarios del río Wierna y el de regantes del río La Caldera--- son de riego (capítulo \ref{cap:amparo}). \textbf{La figura existe y la usan
los que pueden constituirla.}

Y abre una pregunta que este libro no puede contestar: si el Código permite que un conjunto de
vecinos administre el agua de una toma común, \textbf{por qué no hay consorcios de abastecimiento en los parajes
que el capítulo \ref{cap:redes} muestra sin red}. La respuesta puede ser tan simple como que
constituir uno requiere concesiones previas, o tan compleja como que nadie se los propuso.
\end{cautela}

\textbf{Y ahora sí hay medición.}

\begin{ficha}{Censo Nacional 2022, hogares sin servicio de red, por fracción censal}
\textbf{Sin agua por red pública:} fracción 66077-01, \textbf{32,5\,\%}; fracción 66077-02,
\textbf{17,6\,\%}; fracción 66077-03, \textbf{58,0\,\%}.

\textbf{Con el desagüe del inodoro no conectado a red pública:} \textbf{96,3\,\%}, \textbf{95,7\,\%}
y \textbf{95,7\,\%} respectivamente.

En la Capital, sobre cincuenta fracciones, el promedio es de \textbf{5,8\,\%} sin agua de red y
\textbf{13,6\,\%} sin cloaca.
\end{ficha}

Los cuadros publicados del INDEC confirman el orden a escala del departamento, contando personas y no hogares. De las 12.246 personas en viviendas particulares, \textbf{3.407 ---el 27,8 por ciento--- no reciben agua de red pública}: 1.395 se abastecen de perforaciones o pozos, \textbf{1.333 dependen de transporte por cisterna, agua de lluvia, río, canal, arroyo o acequia}, y 679 de otras fuentes. Y sólo 419 personas ---el 3,4 por ciento--- descargan el inodoro a red cloacal: el 96,6 por ciento restante lo hace a cámara séptica, pozo ciego u hoyo, o no tiene baño \textit{(INDEC, Censo 2022, resultados definitivos, cuadros 2.17.11 y 3.17.11 de condiciones habitacionales de la población, Salta)}.


\textbf{Entre uno de cada seis y casi seis de cada diez hogares del departamento no tienen agua
de red, y prácticamente ninguno tiene cloaca.}

Las tres fracciones de La Caldera figuran entre las treinta peores de las ciento sesenta de la
provincia en conexión cloacal. Y la fracción \textbf{66077-03} ---la misma que tiene el
cuarenta y tres por ciento de sus viviendas sin residentes permanentes, según el capítulo
\ref{cap:poblacion}--- es la que peor está en agua: \textbf{puesto 18 de 160} \textit{(el procesamiento propio del capítulo \ref{cap:poblacion} registró 159 fracciones; la diferencia de una no se pudo explicar, de modo que el puesto debe leerse como aproximado)}.

\begin{cautela}
\textbf{Esto es lo que este capítulo venía afirmando sin poder medirlo, y conviene registrar el
cambio.} Hasta la obtención de este dato, la frase «el departamento no tiene red» era una
síntesis de expedientes y obras proyectadas. Ahora es una proporción, con su fuente y su fecha.

\textbf{Lo que el censo mide es el hogar, no la obra.} Dice cuántos hogares declararon no tener
conexión en 2022; no dice si hay red y no están conectados, ni qué obras se hicieron después. La
ampliación de red del barrio Santiago Apóstol que el capítulo \ref{cap:redes} registra es de
septiembre de 2026 y no está en estos números.

\textbf{Y hay una lectura que este libro no hace.} Un porcentaje alto de hogares sin cloaca en
un departamento de baja densidad, con lotes grandes y pozo absorbente, no equivale sin más a una
carencia: puede ser la solución técnica razonable para esa densidad. \textbf{Lo que sí dice el
número es que la red no está}, y que quien vive ahí resuelve por su cuenta lo que en la capital
resuelve el Estado.
\end{cautela}

\section{1944: conceder sin saber}

Cincuenta y cinco años antes del Código, la Provincia ya concedía agua de esta cuenca
declarando por escrito que no sabía si el agua estaba.

\begin{ficha}{Decreto Nº 4614-H --- Salta, 20 de setiembre de 1944 --- Expte.\ Nº 17.870/1944}
A pedido del Sub-Jefe Divisional Tracción Norte de la Administración General de los
Ferrocarriles del Estado ---que pide tomar «un caudal de cinco litros por segundo con un
total diario de 100.000 litros»--- y con informe de la Dirección General de Hidráulica, el
artículo 1º concede autorización para derivar del \textbf{río Mojotoro} un caudal de
\textbf{cinco litros por segundo}, para proveer de agua a las locomotoras que circulan entre
Güemes y Salta. \textit{Las dos cifras del pedido no concuerdan ---cinco litros por segundo
durante un día son 432.000 litros; 100.000 litros equivalen a ese caudal durante unas cinco
horas y media---: así en el original, y el decreto no lo aclara.}

\textbf{Art.\ 2º.} «\textit{El agua se tomará en el sitio y en la forma indicada en el plano
adjunto.}»

\textbf{Art.\ 3º.} «\textit{Esta autorización se concede con carácter precario \textbf{hasta
tanto entre en vigencia el decreto ley de agua de la Provincia} y siempre que ello no perjudique
a terceros.}»

\textbf{Art.\ 4º.} «\textit{Déjase expresamente salvada la responsabilidad de los poderes
públicos de la Provincia \textbf{en cuanto a la existencia en las diferentes épocas del año del
caudal de agua solicitado y autorizado}.}»

Firman Arturo S.\ Fassio, Interventor Federal, e Ismael Casaux Alsina. B.O.\ Nº 2.104,
22 de septiembre de 1944, página 6.
\end{ficha}

Los tres artículos dicen, cada uno, algo que este libro venía sosteniendo por otras vías.

El artículo 3º dice que en septiembre de 1944 \textbf{Salta no tenía régimen legal de aguas en
vigencia}, y que el Ejecutivo lo sabía y concedía igual, en precario, contra una ley futura.

El artículo 4º dice que el organismo que concedió cinco litros por segundo \textbf{no estaba en
condiciones de afirmar que el río los llevara todo el año}, y se exoneró por anticipado de las
consecuencias. Es la misma ausencia de medición que el apartado siguiente encuentra en el
Código de Aguas ---sancionado en 1998, vigente desde 1999---, escrita con todas las letras en 1944: se concede un caudal y se declara no
saber si existe.

Y el artículo 2º remite la ubicación de la toma «al plano adjunto», \textbf{que no se publica}.
Sin ese plano no puede establecerse en qué punto del curso quedó emplazada la obra ---ni
siquiera en qué departamento---, quiénes quedaban aguas abajo ni si hubo oposición. Es, en
1944, exactamente la estructura de los anexos de línea de ribera de 2014 que describe el
capítulo \ref{cap:ribera}: el acto se publica, la ubicación se remite a una pieza que no.

\begin{cautela}
El decreto no nombra departamento. Se incorpora aquí porque el río Mojotoro es el colector de
esta cuenca y da nombre a la localidad que otros actos del mismo año sitúan en el departamento
La Caldera, no porque el Boletín lo diga. \pendiente{plano adjunto al expediente 17.870/1944,
Dirección General de Hidráulica, para ubicar la toma}
\end{cautela}

\section{Y el Código mandó medirla, en 1999}

Hay una obligación que el Código de Aguas impuso y que este libro no puede dar por cumplida.

\begin{ficha}{Ley 7.017, artículos 46, 258 y 259 --- texto literal}
\textbf{Art.\ 46, último párrafo.} «\textit{A partir de la vigencia de este Código \textbf{no se
otorgarán concesiones de ejercicio permanente mientras no sean aforadas sus fuentes de
provisión}.}»

\textbf{Art.\ 258.} «\textit{Durante cinco (5) años a partir de la vigencia de este Código, la
Autoridad de Aplicación \textbf{deberá establecer por fuente, el aforo del agua pública}, cuyo
promedio será denominado ``De Partida''.}» Una vez determinado, se renueva cada diez años.

\textbf{Art.\ 259.} «\textit{Sólo se podrán otorgar concesiones con dotación permanente de agua
únicamente después de determinado el aforo definitivo y previo reajuste de las dotaciones
existentes.}»

Y el \textbf{artículo 161} manda a la Autoridad, en colaboración con la autoridad sanitaria,
hacer un \textbf{inventario de las aguas} dentro de los cinco años de promulgado el Código,
registrarlo y \textbf{actualizarlo anualmente}.
\end{ficha}

El Código se publicó en enero de 1999 y entró en vigencia noventa días después. Los cinco años
de los artículos 258 y 161 vencieron, por lo tanto, en \textbf{2004}.

\begin{cautela}
\textbf{Este libro no afirma que esas obligaciones se hayan incumplido.} No consultó el Libro
de Catastro, no tiene la serie de aforos de partida ni el inventario de aguas, y es
perfectamente posible que existan y no sean públicos.

Lo que sí registra es que \textbf{buscó y no los encontró}, y que la propia Secretaría, en la
presentación que el capítulo \ref{cap:ribera} transcribe, describe una prelación de métodos
cuyo primer escalón ---aforos con estadística no menor de cincuenta años--- supone
exactamente esos registros.

\textbf{Y registra la consecuencia normativa}, que no depende de ninguna verificación: si una
fuente no está aforada, el artículo 46 prohíbe otorgar sobre ella concesiones de ejercicio
permanente. \textbf{La nómina de concesiones permanentes sobre los cursos de este departamento,
puesta al lado de la nómina de fuentes aforadas, contesta sola una pregunta que este libro dejó
abierta.} Las dos nóminas están pedidas en el apéndice \ref{cap:pedidos}.
\end{cautela}

\section{Lo que la obra se llevó}

Hay un costo de las obras de captación que no figura en ningún expediente de expropiación y
que el diagnóstico social del Plan de Manejo recogió de boca de los vecinos.

Durante la construcción del acueducto de Campo Alegre \textbf{se eliminó el bosque ribereño}.
Lo dice el relevamiento y lo dicen los testimonios que el capítulo \ref{cap:expropiacion}
transcribe: «eso cuidaba para que el agua no pasara, pero ahora se lleva puesto todo lo que
encuentra», «culpa de eso el río casi se lleva el cementerio».

\textbf{La obra que lleva agua a la capital se llevó la defensa natural del pueblo.} Y el
departamento viene construyendo defensas ---de rama y piedra primero, de hormigón desde 1980--- desde 1910, según el capítulo
\ref{cap:defensas}, sin que ninguna de las dos series ---la de captación y la de protección---
se haya leído nunca junto a la otra.

Y hay una norma que conviene poner al lado de ese testimonio.

\begin{ficha}{Ley 7.017, artículo 170 --- último párrafo}
«\textit{En todos los casos \textbf{para la tala de árboles situados en las márgenes de los
cursos o depósitos de aguas naturales o artificiales se requerirá permiso de la Autoridad de
Aplicación}.}»

El mismo artículo faculta a la Autoridad a fijar áreas de protección de cuencas donde no se
permita el pastaje, la tala ni la alteración de la vegetación, y a \textbf{disponer la
plantación de árboles o de bosques protectores}.
\end{ficha}

\begin{cautela}
No se sigue de esto ninguna irregularidad. El permiso pudo haberse otorgado, y una obra
provincial de captación tramitada por la propia Autoridad de Aplicación difícilmente carezca
de los actos que ella misma debe dictar.

Lo que este libro registra es que \textbf{la facultad de exigir reposición existía y existe},
en el mismo artículo: la Autoridad puede disponer la plantación de bosques protectores. El
reclamo vecinal por inventarios forestales en la costanera, que el capítulo
\ref{cap:expropiacion} recoge, tiene entonces norma donde apoyarse.
\end{cautela}

\section{Quién decide sobre esta agua}

Hay una tercera línea, más lenta que las obras, que el libro documenta sin nombrarla: el agua
de este departamento pasó de decidirse acá a decidirse en la capital.

En \textbf{1885} el municipio concedía agua. En \textbf{1914} una ordenanza municipal
reglamentaba la distribución de la totalidad del caudal del río Santa Rufina y reconocía
derechos por uso tradicional; el capítulo \ref{cap:expropiacion} muestra que esa ordenanza
llegó a fundar la denegación de una concesión provincial. En \textbf{1942} un derecho otorgado
por el municipio en 1885 se inscribe en los registros provinciales.

Y en \textbf{1951} el movimiento se vuelve masivo: siete actos en un solo año sobre derechos
de agua del departamento, casi todos \textbf{reconocimientos} de derechos ya ejercidos que
pasan a constar en un registro provincial, con expediente, edicto y decreto del gobernador.
El capítulo \ref{cap:redes} lo consigna y señala el dato que importa: \textbf{el municipio no
interviene en ninguno de los siete}.

\begin{cautela}
Vale aquí la cautela del capítulo \ref{cap:redes} sobre 1951: no está probado que ése sea el año en que la Provincia completó el traspaso.

Lo que sí puede sostenerse es la forma del movimiento: \textbf{un derecho que se ejercía sin
registro pasa a existir en un registro que no lleva el municipio}. No se le quita el agua a
nadie. Se cambia quién anota.
\end{cautela}

\section{Y las concesiones sí se publican, una por una}

Este libro no encontró la nómina de concesiones de agua pública del departamento (capítulo \ref{cap:redes}).
La nómina no se publica como documento, pero \textbf{las concesiones se publican de a una}, por
cinco días cada una, bajo el tipo de aviso «concesiones de agua pública». El listado oficial
devuelve, para este departamento, una serie que va de 2004 a 2025.

\begin{ficha}{Algunas de las gestiones publicadas}
\textbf{2004.} Francisco José y Monserrat Alicia Ibáñez Figuerola, expedientes 2.253/56 y 34-3.681/03 (primera publicación del expediente de 1956).

\textbf{2005.} \textbf{Finca La Calderilla}, expediente \textbf{34-7.229/48} y otro ---un expediente de 1948---.

\textbf{2007.} Expediente 34-3.658/03, a nombre de una particular.

\textbf{2007--2011.} Maximiliano Rodrigo Jiménez, \textbf{arroyo Las Vertientes}, expediente
0050034-8.402/2007.

\textbf{2009.} Da Souza Correa Hunt Lubel, expediente 34-189.304/99 ---un expediente iniciado
diez años antes---.

\textbf{2010.} \textbf{El Gallinato S.R.L.}, expediente 0050034-15.196/10.

\textbf{2010.} José y Montserrat Ibáñez Figuerola, \textbf{finca «Vaqueros»}, expediente
\textbf{34-2.253/56} ---un expediente de 1956 que se publica cincuenta y cuatro años después---.

\textbf{2011.} Magdalena Vidal\index{Vidal, Magdalena}, expediente 0050034-23.217/10.

\textbf{2012.} Catastro 4068, expediente 0090034-53.406/12. \textbf{2013.} Matrícula 3721, ríos Wierna y Bermejo, expediente 0090034-94.536/13, que vuelve a publicarse en 2018. \textbf{2014.} Catastro 1700 (expte.\ 0090034-25.217/2014), catastro 2907 (0090034-168.839/14) y Finca Lesser (34-6.147/05).

\textbf{2015.} Río La Caldera (34-211.560/2014, el dren); expedientes 0090034-159.566/2012 (con dos agregados) y 0090034-110.501/2012. \textbf{2016.} Consorcio Valle Alegre (0090034-86.490/2013). \textbf{2017.} \textbf{Las Vertientes--Eco Pueblo}, río Wierna (0090034-263.317/2015; capítulo \ref{cap:redes}).

\textbf{2018--2019.} Expedientes 0090034-254.833/17, 0090034-179.036/2018 y 0090034-88.064/2015-0/\textbf{14}. \textbf{2021--2022.} 0090034-138.006/2020, 0090034-81.199/2018 y 0090034-206.106/2022. (Dos avisos de 2022 que el buscador devuelve para «La Caldera» ---expedientes 0090034-172.447/2012 y 0090034-32.055/2022--- son del departamento Rosario de Lerma, río Toro: el texto escribe «Dpto.\ Rosario de Lerma, Caldera».) \textbf{2023--2025.} 0090034-121.463/2022, 0090034-235.963/2022 y 0090034-189.047/2024; ninguno de estos títulos nombra el curso ni el peticionante.

\textbf{2025.} El Chalchal S.R.L.\ y otro, expediente 0090034-226732/2024-0: como cotitular del catastro 5925, gestiona agua del \textbf{arroyo Los Mártires} para \textbf{uso pecuario de cien cabezas}, 6,5 m³ por día, con carácter permanente (edicto del 24 de abril de 2025; cinco publicaciones en mayo, por \$35.062,50); y expediente
0090034-66789/2025-0 sobre el río La Caldera: el titular registral del catastro 3152 gestiona riego para \textbf{2,5 hectáreas}, con carácter \textbf{eventual} y una dotación de 1,3125 litros por segundo, «conducidas mediante acequia comunera del sistema correspondiente al \textbf{Consorcio del río Mojotoro}». Edicto del 14 de noviembre de 2025; cinco publicaciones en diciembre, por \$36.125.
\end{ficha}

\begin{ficha}{Lo que dicen los edictos abiertos}
\textbf{Agua para loteos} (todas de carácter permanente):
\begin{itemize}[leftmargin=*,nosep]
\item \textbf{Cerro de Buena Vista}, catastros 1713 y 1905 (expte.\ 0090034-159.566/2012; edicto del 30/07/2015): dren horizontal en el río Vaqueros, 63,75 m³/día para 51 lotes, y riego eventual de 3,81 ha de espacios comunes.
\item \textbf{La Misión Etapa III}, catastro 3616, El Chalchal S.R.L.\ (expte.\ 0090034-110.501/2012; edicto del 06/10/2015): pozo perforado, 9 m³/hora durante 22 horas, para 580 habitantes en 116 lotes.
\item \textbf{Las Vertientes--Eco Pueblo}, catastro 4095 (expte.\ 0090034-263.317/2015; edicto del 07/02/2017): acequia El Alto del río Wierna, 600,48 m³/día para 630 familias.
\item \textbf{Valle Hermoso}, catastro 3307, Valle Hermoso S.R.L.\ (expte.\ 0090034-88.064/2015-0/14; edicto del 10/04/2019): dren horizontal, 15,71 m³/hora durante 16 horas (0,092 hm³/año).
\end{itemize}
Las de Cerro de Buena Vista, La Misión y Valle Hermoso se publican como «finalización del trámite» de la concesión.

\textbf{Riego que sigue a la subdivisión:}
\begin{itemize}[leftmargin=*,nosep]
\item 2005 (exptes.\ 34-7.229/48 y 34-4.604/04): división de la concesión de la finca La Calderilla, catastro de origen 38, reconocida por el Decreto 9838, por subdivisión según el plano 336; a la matrícula 3127 le corresponden 4,69 ha permanentes y 2,462 l/s del río La Caldera. Firma todavía la Agencia de Recursos Hídricos.
\item 2018 (expte.\ 0090034-94.536/13): división de la concesión del catastro de origen 3721 ---Finca Wierna, el inmueble de las Resoluciones 74/15 y 293/15--- a favor de la matrícula 5009: 18 ha permanentes y 9,45 l/s del río Wierna, margen izquierda, «cuenca del río Bermejo». Es el mismo expediente publicado en 2013. En 2022 (expte.\ 0090034-81.199/2018) otra fracción del mismo origen, el catastro 5010, recibe 15,33 ha permanentes y 8,05 l/s: el riego de Finca Wierna se reparte, al menos, entre dos matrículas nuevas.
\item 2021 (expte.\ 0090034-138.006/2020): el riego del catastro 2907 ---10 ha permanentes, 5,25 l/s del río Wierna--- pasa al catastro 6582. En 2014 el mismo riego había pasado del catastro 126 al 2907.
\item 2022 (expte.\ 0090034-206.106/2022): subdivisión proporcional del catastro 169 en los 3436 y 3437, con 60 ha eventuales y 31,5 l/s cada uno, del río Las Pavas.
\item 2023 (expte.\ 0090034-235.963/2022): anulación y subdivisión para anexión de los catastros 84, 2264 y 2932 ---8 ha permanentes y 4 eventuales--- en los 7155 y 7156, por la acequia La Calderilla. Las permanentes se reparten completas (1 + 7); de las eventuales se asignan 3,5 de las 4.
\end{itemize}

\textbf{Y los siete edictos tienen un número en común.} En todos, sin excepción, la dotación
resulta de multiplicar la superficie por \textbf{0,525 litros por segundo y por hectárea}:
4,69 ha y 2,462 l/s en 2005; 18 y 9,45 en 2018; 10 y 5,25 en 2021; 60 y 31,5 y 15,33 y 8,05
en 2022; 2,5 y 1,3125 en 2025; y los 10 ha y 5,25 l/s del catastro 2907 que el capítulo
\ref{cap:redes} transcribe para 2014. Tres cursos distintos ---el río La Caldera, el Wierna
y el Las Pavas---, veinte años de diferencia, concesiones permanentes
y eventuales: el mismo coeficiente, con las dotaciones redondeadas a dos o tres decimales. Y no empieza con estos edictos: los dos reconocimientos de 1948 que transcribe el capítulo \ref{cap:siglo} ---La Helvecia y El Angosto--- dan la misma razón.

\begin{cautela}
\textbf{Eso no es una irregularidad y puede ser lo correcto.} Un coeficiente uniforme de
dotación por hectárea es una técnica administrativa razonable, probablemente fijada en una
resolución que este libro no leyó, y evita la discrecionalidad caso por caso.

\textbf{Lo que muestra es cómo se asigna el agua en esta cuenca.} La dotación no surge de
medir el curso: surge de multiplicar la superficie por una constante. Puesto al lado de los
artículos 46, 258 y 259 ---que prohíben otorgar dotación permanente sin aforo previo y mandaban
tener el aforo de partida en 2004---, el dato cierra el apartado anterior por otro camino:
\textbf{el Estado reparte agua por fórmula porque no tiene con qué repartirla por medición}.
Cuál es esa resolución, desde cuándo rige ---al menos desde 1948--- y sobre qué base se calculó el 0,525 es un pedido de
una línea, y queda sumado al apéndice \ref{cap:pedidos}.
\end{cautela}

\textbf{Abastecimiento fuera de los loteos}: en 2018, un grupo de vecinos gestiona 10 m³/día permanentes de vertientes «en el catastro de mayor extensión Nº 4.097» para cuarenta personas (expte.\ 0090034-254.833/17). Es la matrícula sobre la que la Provincia fijó líneas de ribera en 2016, 2021 y 2022 (capítulo \ref{cap:ribera}).

\textbf{Y usos nuevos}: una concesión de energía hidráulica al pie del dique Campo Alegre, en 2019 (capítulo \ref{cap:expropiacion}); en 2023, riego eventual para un vivero comercial de 8.000 m² con 0,40 l/s del río Wierna (expte.\ 0090034-121.463/2022); y en 2024 el \textbf{Arzobispado de Salta}, cotitular del catastro 1673, gestiona riego eventual para 3 ha (1,575 l/s del río La Caldera, margen izquierda) y abrevado de cincuenta cabras, con la previsión de captar «en épocas críticas» de una vertiente innominada que nace en la propiedad (expte.\ 0090034-189.047/2024; edicto del 11/12/2024).
\end{ficha}

\begin{cautela}
Dos cosas se siguen de estos edictos. La primera es que los loteos del corredor tienen su agua documentada, y de cuatro fuentes distintas: un pozo, dos drenes y una acequia. Las cuatro concesiones son permanentes, y la de La Misión, captada de un pozo, lo es pese a que el artículo 143 declara eventuales las subterráneas (capítulo \ref{cap:loteo}). Dos de los cuatro inmuebles ---el 3616 y el 1905--- son los mismos sobre los que la Provincia fijó líneas de ribera en 2014, y el club Las Vertientes tuvo la suya ese año sobre la matrícula de su plano. Ribera y agua se tramitaron para los mismos predios, en los mismos años, ante el mismo organismo.

La segunda es que el padrón de riego sí se divide cuando la tierra se divide: hay un procedimiento, se publica y se usa desde 2005. En uno de los casos de 2023, sin embargo, la suma de lo asignado no alcanza lo que tenían los catastros de origen. Y ninguno de los edictos dice si el derecho se otorgó.
\end{cautela}

\begin{cautela}
\textbf{Esto es un camino, no un hallazgo.} Recorrido entero, el listado suma unas treinta gestiones propias del departamento entre 2004 y 2025, publicadas cinco días cada una, además de avisos de línea de ribera clasificados bajo el mismo tipo. Los títulos dan el peticionante, el expediente y a
veces el curso; \textbf{no dan el caudal, ni el uso, ni si la concesión se otorgó}. Cada uno
requiere abrir su edicto.

Lo que sí establece la serie es que \textbf{la información existe y está publicada}: quien
quiera armar la nómina de concesiones de esta cuenca puede hacerlo edicto por edicto, sin pedirle
nada a nadie. Este libro no lo hizo y lo deja señalado.

\textbf{Y trae tres datos que valen por sí solos.} El arroyo que da nombre al loteo Las
Vertientes tuvo una gestión de concesión a nombre de un particular entre 2007 y 2011. La
sociedad El Gallinato S.R.L.\ gestionó agua pública en 2010, cinco años antes de que se
aprobara sobre su matrícula el club de campo que el capítulo \ref{cap:loteo} documenta. Y un
expediente iniciado en \textbf{1956} seguía publicándose en 2010, lo que da la medida de los
tiempos de este trámite.
\end{cautela}

\section{El agua que no se registra}

Y hay una cuarta categoría, que el capítulo \ref{cap:redes} encuentra en un decreto de 1939.
Un arroyo que \textbf{nace y muere dentro de una misma finca} no se inscribe en el registro de
aguas públicas, porque el uso de un curso de dominio privado no es una concesión.

\textbf{Sobre esa agua el Estado no lleva registro, y por lo tanto este libro tampoco puede
saber cuánta es.} En un departamento de cabecera de cuenca, con decenas de cursos menores que
nacen y mueren dentro de fincas grandes, la categoría no es marginal. Es, literalmente, agua
invisible.

\section{Lo que pasa mientras tanto en el cauce}

Dos series más, que los capítulos \ref{cap:ribera} y \ref{cap:resistencias} documentan y que
solo puestas juntas dan la escena completa.

\textbf{Veinticinco actos de línea de ribera} entre 2013 y 2015, con la asimetría que
el capítulo \ref{cap:ribera} reconstruye: la misma matrícula, el mismo organismo, cinco meses
de diferencia, y una publica ciento trece pares de coordenadas mientras la otra las remite a
la sede.

\textbf{Veinte concesiones de extracción de áridos} sobre las márgenes, de las cuales cinco
tienen la evaluación ambiental aprobada y seis constan expresamente como no aprobadas. Las veinte, vigentes según la nota de 2022; el informe de Minería de 2023 registra, sin embargo, la renovación de Los Yacones rechazada en 2020 (capítulo \ref{cap:amparo}). Y dentro de la \textbf{Zona de Exclusión de Aguas del Norte} existía un
convenio que a 2022 ya no estaba vigente, razón por la cual ---dice la autoridad minera--- se paralizaron las actividades extractivas; el acuerdo que lo reemplazó se formalizó ese mismo año por la Resolución SMyE 20/2022, que no se publicó (capítulo \ref{cap:amparo}).

\textbf{La misma empresa que lleva el agua a la capital tenía un convenio para ordenar los
laboreos sobre el cauce del que la saca.} Cuando el convenio cayó, lo reemplazó una resolución que no se publicó.

\section{Lo que no existe}

Puesto todo junto, lo que falta se nombra solo.

Hay una cuenca. Hay dos municipios que la comparten y no administran el conjunto del cauce.
Hay una empresa provincial que capta, una secretaría que determina riberas, otra que concesiona
áridos, una tercera que evalúa impacto, un juzgado que aprueba un plan y una ley de 1929 que le
dio potestad de policía a un municipio de otro departamento.

\textbf{No hay autoridad de cuenca.} La única figura metropolitana con ley, organismo y presupuesto es la del transporte de pasajeros; el Corredor Intermunicipal de 2009 existe como convenio, sin actos de funcionamiento conocidos, y el Comité Metropolitano de Ordenamiento Territorial que invoca la Ordenanza 535/18 de Vaqueros acordó criterios para urbanizar, no para administrar el río (capítulo \ref{cap:redes}).

\section{Qué afirma este capítulo y qué no}

\begin{cautela}
\textbf{Afirma} que las obras de captación de este siglo y medio llevan el agua en una sola
dirección; que el registro de los derechos se trasladó del municipio a la provincia; que el
cauce está simultáneamente concesionado para extraer, delimitado de manera desigual y ocupado
por loteos; y que ningún organismo administra el conjunto.

\textbf{No afirma} que exista un plan detrás de eso. No hay un acto que ordene la secuencia, y
este relevamiento no encontró ninguno. Cada pieza tiene su expediente, su fundamento y su firma
---y varias tienen buenos fundamentos---.

\textbf{Y no afirma que la dirección sea reversible.} Una ciudad de más de medio millón de
habitantes necesita el agua que baja de esta cuenca, y ningún ordenamiento razonable
propondría que deje de bajar. La pregunta que el libro deja abierta no es si el agua debe
salir, sino \textbf{qué le corresponde al territorio que la produce}, y quién está obligado a
responder esa pregunta.

Hoy, nadie.
\end{cautela}

